\subsection{Permute \& Reject}

  \begin{figure}[h]
    \begin{center}
      \includegraphics[width=0.3\textwidth]{img/permute16.eps}
    \end{center}
    \caption{Polygon mit 16 Punkten (Permute \& Reject)}
    \label{fig:permute16}
  \end{figure}

  \emph{Permute \& Reject} ist der einfachste in der Software enthaltenene
  Algorithmus zur Erzeugung von simplen Polygonen, implementiert nach einer
  Beschreibung von Martin Held in~\cite{held98polygons}.

  \subsubsection{Algorithmus}

    Der Algorithmus basiert auf der zufälligen Erzeugung von $n$ Kanten zwischen
    $n$ gegebenen Punkten, sodass von jedem Punkt zwei Kanten ausgehen.
    Anschließend wird überprüft, ob das entstandene Polygon \enquote{simpel},
    d.h. überschneidungsfrei ist. Falls nicht, wird das Polygon verworfen und
    von vorne begonnen.

\begin{code}[caption={Permute \& Reject},label=permutelisting,mathescape=true]
retry:
$P$ $\leftarrow$ given set of points
$p$ $\leftarrow$ remove random point from $P$
while ($P$ not empty)
  $q$ $\leftarrow$ remove random point from $P$
  draw edge from $p$ to $q$
  $p$ $\leftarrow$ $q$
if resulting polygon is NOT simple
  goto retry
\end{code}

  \subsubsection{Eigenschaften}

    Ausgehend von einer beliebigen Punktmenge kann Permute \& Reject sämtliche
    möglichen Polygone erzeugen. Jedoch lässt die enorme Laufzeit des
    Algorithmus eine praktische Verwendung für Polygone mit vielen Punkten
    nicht zu. Die Laufzeit ist direkt abhängig von der Anzahl der benötigten 
    \enquote{Versuche}, d.h. der Anzahl der erzeugten Permutationen und der
    Tests auf Schnittfreiheit der resultierenden Polygone. Da die Implementierung
    keinerlei \emph{blacklisting} betreibt, ist es möglich, dass bereits
    getestete Permutationen erneut erzeugt werden, sodass die Laufzeit theoretisch
    nicht beschränkt ist. In unseren Experimenten konnten wir kein Polygon mit mehr
    als 23 Punkten erzeugen.

    Die erwartete Laufzeit des Algorithmus ist abhängig von der Anzahl der \enquote{gutartigen} 
    Permutationen in der Menge aller Permutationen der Punktmenge ($n!$ Permutationen). 
    Im \emph{worst case} liegen alle Punkte der gegebenen Punktmenge
    auf der konvexen Hülle der Menge (bspw. bilden die Punkte eine nahezu regelmäßiges
    Vieleck). In diesem Fall gibt es nur eine gültige geometrische Lösung, die allerdings
    durch $2n$ Permutationen ausgedrückt werden kann, da jeder Punkt Startpunkt der Polygonkette
    sein kann und jede gültige Polygonkette auch eine exakte Umkehrung besitzt, die ebenfalls
    das selbe einfache Polygon beschreibt. Für den \emph{worst case} lässt sich hieraus 
    folglich eine erwartete Anzahl von $\frac{n!}{2n}$ Iterationen für den Algorithmus ableiten.

  \subsubsection{Enumerating Permute \& Reject}

    Um die Laufzeit des Algorithmus auch tatsächlich zu beschränken, implementierten wir
    eine zweite Version, die auf einer einfachen Aufzählung aller möglichen Permutationen
    der Punktmenge beruht. Dieser neue, \emph{Enumerating Permute \& Reject} genannte Algorithmus
    benötigt im oben beschriebenen schlechtesten Fall nicht mehr als $\frac{n!}{2n}$ Iterationen.
    Zu unserem Bedauern mussten wir jedoch feststellen, dass der Algorithmus ab einer gewissen
    Punktanzahl in der Regel deutlich größere Laufzeiten aufweist als \emph{Permute \& Reject}.
    Wir vermuten, dass dies aus der Wahl der Steinhaus-Johnson-Trotter-Algorithmus zur iterativen 
    Erzeugung der Permutationen resultiert, da die Permutationen hier stets nur an einer Stelle
    verändert werden.
  